\section{Summary}

\vspace{1cm}

\subsection{Introduction}

This experiment aims to evaluate the performance of compressed sensing using two models on the MNIST dataset: a Variational Auto-encoder (VAE) and a Generative Adversarial Network (GAN). The focus is on image reconstruction from compressed measurements. The MNIST dataset consists of 60,000 training images and 10,000 test images of handwritten digits, each of size 28x28 pixels.

\subsection{Models and methodology}

We employed two approaches to reconstruct the images from compressed measurements:
\begin{enumerate}
    \item \textbf{Variational Auto-encoder (VAE):} This model is designed to generate images similar to those in the training set. The VAE compresses the input image into a lower-dimensional latent space and then reconstructs the image from this representation.
    \item \textbf{Generative Adversarial Network (GAN):} This model is composed of two competing networks, a generative one and a discriminating one. They are trained to minimize their losses:
    \begin{itemize}
        \item The generator tries to generate images from a random latent vector.
        \item The discriminator tries to discern between real images and generated ones.
    \end{itemize}
\end{enumerate}
Each of the two was trained with a latent vector of dimension 20 and of dimension 30. We also trained the fully connected auto-encoder used by Bora et al., so that our results can be compared directly with the ones they publish.
\\
For all the models, we used a random Gaussian matrix for the measurements, and LASSO was run in two sparsifying bases, the pixel basis used by the reference paper on MNIST and the DCT basis. The compression levels varied from 10 to 750 measurements per image, out of the 784 pixels of a full image, allowing us to observe the reconstruction quality across different levels of compression.

\subsection{Algorithm explanation}

This section will explain the algorithm utilizing a generative model for compressed sensing.
\\
Enhancements include a thousand gradient steps on the latent vector and 10 random restarts to improve convergence and error minimization.



\subsection{The MNIST Dataset}

The MNIST (Modified National Institute of Standards and Technology) dataset is a large collection of handwritten digits commonly used for training and evaluating image processing systems. It contains 60,000 training images and 10,000 test images, each being a gray scale image of size \( 28 \times 28 \) pixels. Each image is labeled with a digit from 0 to 9, providing a diverse set of handwriting styles for developing robust digit recognition systems.
\begin{itemize}
    \item \textbf{Size:} 60,000 training images and 10,000 test images.
    \item \textbf{Resolution:} Each image is \( 28 \times 28 \) pixels, resulting in 784 features.
    \item \textbf{Classes:} 10 classes, representing digits 0 through 9.
    \item \textbf{Grayscale:} Pixel values range from 0 (black) to 255 (white).
\end{itemize}
The MNIST images were vectorized, resulting in 784-dimensional input vectors. Each pixel value was normalized to the range $[0, 1]$.



\subsection{Results}

The reconstruction performance was evaluated as the squared error between the original and the reconstructed image, divided by the number of pixels, averaged over ten test images, one for each digit class. The following observations were made:

\begin{enumerate}
    \item \textbf{Performance with few measurements:} the generative priors are far more accurate than either LASSO baseline when measurements are scarce, the VAE priors by the widest margin. Every VAE prior reaches with 75 measurements the accuracy either baseline obtains with 400, a saving of a factor 5.3, within the range of 5 to 10 the authors report.
    \item \textbf{Saturation:} the accuracy of the generative priors stops improving after roughly 200 measurements, because the reconstruction is confined to the range of the generator. From 500 measurements onwards both LASSO baselines become more accurate, as the reference paper predicts, and with 750 measurements the one in the pixel basis recovers the digits almost exactly.
    \item \textbf{Comparison between the models:} the three VAE priors cannot be told apart with ten test images, no comparison among them being significant at any budget. The DCGAN with latent dimension 20 is significantly less accurate than every VAE at nine budgets out of ten; the one with latent dimension 30 also has a higher error at every budget, but ten test images establish it only where measurements are scarce. Both cost more than ninety times as much to invert as the cheapest VAE.
\end{enumerate}

\subsection{Conclusion}

The experiment demonstrated the efficacy of using deep learning models, particularly VAEs, for image reconstruction from compressed measurements. These models significantly reduce the reconstruction error compared to traditional sparse recovery methods when the number of measurements is small, while the traditional method remains preferable when measurements are abundant. The findings highlight the potential of integrating compressed sensing techniques with deep learning to enhance image reconstruction quality in resource-constrained environments.
\\
These results provide a promising direction for further research, particularly in optimizing measurement matrices and exploring other deep learning architectures for compressed sensing applications.
